% Jordan分解的极小性：eta=lambda|_[0,1]-lambda|_[2,3]。(a) Jordan部分eta+=lambda|_[0,1]，eta-=lambda|_[2,3]，总质量各为1；(b) 另一种表示alpha=lambda|_[0,3]，beta=lambda|_[1,2]+2lambda|_[2,3]，总质量各为3。块高为长度测度前的系数，1单位=0.45cm；横轴1单位=1.1cm。
\begin{tikzpicture}[
  lab/.style={font=\figurefont\fontsize{8}{10}\selectfont,text=FigureInk},
  note/.style={font=\FigureNoteFont\fontsize{7.5}{9}\selectfont,text=FigureMuted,align=center},
  ax/.style={draw=FigureStroke,line width=.75pt,-{Stealth[length=2.2mm]}},
  posbar/.style={draw=FigureBlue,fill=FigureBlueFill,line width=.7pt},
  negbar/.style={draw=FigureRed,fill=FigureRedFill,line width=.7pt}]
\def\u{1.1}\def\h{0.45}
\begin{scope}
  \draw[posbar] (0,0) rectangle (\u,\h);
  \draw[ax] (-.3,0)--({3*\u+.5},0);
  \node[lab,anchor=west] at ({3*\u+.6},0) {$\eta^+$};
  \foreach \x in {0,1,2,3}{\draw[FigureStroke,line width=.6pt] ({\x*\u},0)--({\x*\u},-.08); \node[lab,below] at ({\x*\u},-.06) {$\x$};}
  \begin{scope}[yshift=-1.6cm]
    \draw[negbar] ({2*\u},0) rectangle ({3*\u},\h);
    \draw[ax] (-.3,0)--({3*\u+.5},0);
    \node[lab,anchor=west] at ({3*\u+.6},0) {$\eta^-$};
    \foreach \x in {0,1,2,3}{\draw[FigureStroke,line width=.6pt] ({\x*\u},0)--({\x*\u},-.08); \node[lab,below] at ({\x*\u},-.06) {$\x$};}
    \node[note] at ({1.5*\u},-.85) {$\eta^+(X)=\eta^-(X)=1$};
    \node[lab] at ({1.5*\u},-1.4) {(a) Jordan分解};
  \end{scope}
\end{scope}
\begin{scope}[xshift=5.4cm]
  \draw[posbar] (0,0) rectangle ({3*\u},\h);
  \draw[ax] (-.3,0)--({3*\u+.5},0);
  \node[lab,anchor=west] at ({3*\u+.6},0) {$\alpha$};
  \foreach \x in {0,1,2,3}{\draw[FigureStroke,line width=.6pt] ({\x*\u},0)--({\x*\u},-.08); \node[lab,below] at ({\x*\u},-.06) {$\x$};}
  \begin{scope}[yshift=-1.6cm]
    \draw[negbar] (\u,0) rectangle ({2*\u},\h);
    \draw[negbar] ({2*\u},0) rectangle ({3*\u},{2*\h});
    \draw[ax] (-.3,0)--({3*\u+.5},0);
    \node[lab,anchor=west] at ({3*\u+.6},0) {$\beta$};
    \foreach \x in {0,1,2,3}{\draw[FigureStroke,line width=.6pt] ({\x*\u},0)--({\x*\u},-.08); \node[lab,below] at ({\x*\u},-.06) {$\x$};}
    \node[note] at ({1.5*\u},-.85) {$\alpha(X)=\beta(X)=3$，在$[1,3]$上重叠};
    \node[lab] at ({1.5*\u},-1.4) {(b) 同一$\eta$的另一种差};
  \end{scope}
\end{scope}
\end{tikzpicture}
